% !TeX encoding = UTF-8
\documentclass[variant=docente,cover=institutional,mode=screen]{ua-engineering-v6-article}
% !TeX encoding = UTF-8
\UASetUniversity{Universidad Austral}
\UASetFaculty{Facultad de Ingeniería}
\UASetDepartment{Departamento de Computación}
\UASetCampus{Pilar}
\UASetCareer{Ingeniería Informática}
\UASetCommission{Comisión A}
\UASetProfessor{Equipo docente}
\UASetSemester{Segundo cuatrimestre 2026}
\UASetCourse{Sistemas Distribuidos}
\UASetDate{2026}


\UASetClassNumber{16}
\UASetClassTitle{Defensa final ante comité}
\UASetDocKind{Guía docente}

\title{Clase 16 --- Guía docente}
\author{\UAProfessor}
\date{\UADate}

\begin{document}
\maketitle

\section{Objetivo pedagógico profundo}

La meta de esta clase es culminar el curso desplazando al estudiante desde el conocimiento fragmentado hacia la capacidad de exposición, defensa y revisión crítica integral de un sistema distribuido.

La idea que debe quedar instalada es:

\begin{center}
\textbf{la madurez en sistemas distribuidos se evidencia cuando el alumno puede convertir mecanismos en una arquitectura defendible bajo objeciones técnicas reales}
\end{center}

\section{Resultados de aprendizaje esperados}

Al finalizar, el estudiante debería poder:

\begin{itemize}
\item estructurar una defensa arquitectónica integral;
\item usar el marco $D=(P,F,G,S,T)$ como guía de argumentación;
\item responder objeciones hostiles sin caer en slogans;
\item reconocer límites y trabajo futuro sin debilitar la propuesta;
\item evaluar técnicamente propuestas de otros con criterio explícito.
\end{itemize}

\section{Estructura sugerida para 90 minutos}

\subsection*{Bloque 1 (15 min): cambiar de rol}
Abrir explicando que esta clase no es otra “unidad temática”, sino el momento en que el estudiante deja de ser receptor de contenidos y pasa a ser defensor técnico.

\paragraph{Objetivo didáctico}
Marcar el cambio de posición epistemológica: ya no “qué aprendimos”, sino “qué podemos defender”.

\subsection*{Bloque 2 (20 min): estructura de la defensa}
Presentar una secuencia clara de defensa:
\begin{itemize}
\item problema;
\item restricciones;
\item fallas;
\item garantías;
\item arquitectura;
\item operación;
\item trade-offs;
\item límites.
\end{itemize}

\paragraph{Objetivo didáctico}
Dar una plantilla conceptual sólida para ordenar exposiciones.

\subsection*{Bloque 3 (20 min): objeciones hostiles}
Trabajar preguntas típicas de jurado:
\begin{itemize}
\item ¿por qué no un monolito?;
\item ¿por qué no consistencia fuerte en todo?;
\item ¿qué pasa bajo partición?;
\item ¿cómo operás esto?;
\item ¿qué parte no está resuelta?
\end{itemize}

\paragraph{Objetivo didáctico}
Entrenar defensa bajo presión sin perder rigor.

\subsection*{Bloque 4 (15 min): revisión crítica}
Discutir cómo una buena defensa incluye autocrítica:
\begin{itemize}
\item límites;
\item partes frágiles;
\item validación pendiente;
\item trabajo futuro.
\end{itemize}

\paragraph{Objetivo didáctico}
Evitar que el alumno piense que defender bien implica fingir perfección.

\subsection*{Bloque 5 (20 min): rúbrica y evaluación}
Presentar una rúbrica clara y usarla sobre un caso o mini-presentación.

\paragraph{Objetivo didáctico}
Transformar la defensa en práctica evaluable y no en intuición difusa.

\section{Qué explicar y por qué}

\subsection*{1. Estructura de defensa}
Porque muchos alumnos saben contenido pero no saben organizarlo como argumento.

\subsection*{2. Marco $D=(P,F,G,S,T)$}
Porque resume la lógica del curso y da una columna vertebral reusable.

\subsection*{3. Objeciones}
Porque obligan a consolidar comprensión real y no recitación.

\subsection*{4. Revisión crítica}
Porque la honestidad técnica es parte esencial de una buena defensa.

\subsection*{5. Rúbrica}
Porque hace explícito qué se considera valioso en una propuesta madura.

\section{Preguntas disparadoras recomendadas}

\begin{itemize}
\item ¿Qué problema concreto justifica esta complejidad?
\item ¿Qué costo estás aceptando y por qué?
\item ¿Qué parte de tu propuesta es más débil?
\item ¿Qué respuesta darías si el jurado no acepta tu supuesto principal?
\item ¿Qué diferencia tu propuesta de una arquitectura “de moda”?
\end{itemize}

\section{Errores frecuentes del alumnado}

\begin{itemize}
\item presentar componentes antes de formular el problema;
\item responder objeciones con slogans;
\item ocultar trade-offs;
\item confundir sofisticación con calidad;
\item no integrar observabilidad, operación e incidentes;
\item no reconocer límites o hipótesis.
\end{itemize}

\section{Ejercicio en vivo sugerido}

Dividir la clase en pequeños grupos. Asignar a cada uno un dominio:
\begin{itemize}
\item e-commerce global;
\item ledger financiero;
\item plataforma de eventos.
\end{itemize}

Cada grupo prepara una mini-defensa en 5 minutos usando:
\[
D = (P,F,G,S,T)
\]

Luego el resto de la clase actúa como comité hostil con 2 o 3 preguntas fuertes.

\paragraph{Objetivo}
Practicar simultáneamente formulación, defensa y evaluación.

\section{Momento crítico de la clase}

El punto más importante ocurre cuando el estudiante comprende que:

\begin{center}
\textbf{la defensa arquitectónica madura no vende una solución perfecta: muestra una solución coherente, limitada y razonada}
\end{center}

Ese momento debe enfatizarse, porque resume el espíritu del cierre del curso.

\section{Criterio de éxito docente}

La clase salió bien si el estudiante puede:

\begin{itemize}
\item estructurar una defensa sin improvisación caótica;
\item responder objeciones con claridad;
\item reconocer límites sin derrumbar su propuesta;
\item evaluar críticamente otra defensa usando una rúbrica explícita.
\end{itemize}

\section{Cierre del curso}

Esta clase debe cerrar el curso con una idea fuerte:

\begin{quote}
Saber sistemas distribuidos no es solo conocer mecanismos, sino poder pensar, diseñar, operar y defender sistemas complejos bajo restricciones reales.
\end{quote}

\begin{uakeybox}
El curso termina cuando el estudiante ya no solo puede explicar qué hace un mecanismo, sino cuándo conviene usarlo, qué costo implica y cómo defender esa decisión ante crítica seria.
\end{uakeybox}


\section{Plan de conducción ampliado}
La clase debe alternar explicación, predicción y contraste. Antes de revelar cada mecanismo, pedir al grupo que anticipe qué puede salir mal; después, volver al mismo escenario y preguntar qué propiedad mejora y qué costo aparece. Cada bloque conceptual debe cerrar con una decisión de diseño observable.
\subsection*{Secuencia recomendada}
\begin{enumerate}
\item Recuperar el prerequisito de la clase anterior.
\item Presentar un caso mínimo que falle y solicitar hipótesis.
\item Formalizar el concepto central de Defensa final ante comité, distinguiendo seguridad, progreso y supuestos.
\item Resolver un ejemplo completo en pizarra, incluyendo una falla intermedia.
\item Comparar una alternativa de diseño y explicitar trade-offs.
\item Cerrar con una pregunta del Trabajo Final y una pregunta de defensa oral.
\end{enumerate}
\section{Diagnóstico de comprensión durante la clase}
No preguntar solamente ``¿se entendió?''. Usar preguntas discriminantes: ``¿qué cambia si el mensaje llega dos veces?'', ``¿qué propiedad sigue siendo cierta si cae un nodo?'', ``¿esto demuestra seguridad o progreso?'', ``¿qué observa un cliente externo?''. Si el estudiante responde solo con nombres de patrones, pedir una ejecución concreta.
\section{Criterios para corregir ejercicios}
Una respuesta excelente declara supuestos, construye una secuencia verificable, nombra la garantía con precisión, reconoce un costo y conecta la decisión con el dominio. Penalizar afirmaciones absolutas sin condiciones.
\section{Vinculación con el Trabajo Final Integrador}
Reservar entre 5 y 10 minutos para que cada grupo registre una decisión con: problema, alternativa elegida, alternativa descartada y razón. Esto crea trazabilidad entre las clases y las entregas acumulativas.
\section{Lectura docente previa}
Revisar la bibliografía indicada en los apuntes y preparar al menos un contraejemplo que no aparezca literalmente en las diapositivas. El objetivo es poder responder preguntas de frontera sin convertir la clase en una lectura del deck.

\section{Hito del Trabajo Final: defensa}
La defensa debe evaluar la integración del sistema y no una enumeración de tecnologías. El comité debe exigir evidencia ejecutable de las decisiones implementadas, formular objeciones sobre fallas y trade-offs, y distinguir con claridad entre lo implementado, lo simulado y lo propuesto arquitectónicamente.

\section{Arquitectura didáctica v9 para 180 minutos}
\subsection*{Bloque 1 - desafío inicial (20 min)}
Presentar una situación donde aparezca presentación descriptiva, respuesta evasiva, límite no reconocido y rúbrica inconsistente. Pedir predicciones individuales antes de introducir vocabulario. El objetivo es hacer visibles los supuestos intuitivos y recuperar el prerequisito de la clase anterior.

\subsection*{Bloque 2 - construcción conceptual (35 min)}
Formalizar presentación ante comité, objeciones, revisión crítica, rúbrica y defensa oral individual. Separar seguridad, progreso y experiencia del usuario. Explicar no solo \emph{qué} hace cada mecanismo, sino \emph{por qué} existe y qué supuesto permite que funcione.

\subsection*{Bloque 3 - modelo y figura (35 min)}
Construir en pizarra una línea temporal, grafo, log, máquina de estados o tabla de decisiones. Recorrer una ejecución nominal y una adversarial. La representación debe quedar como referencia para los apuntes y el ejercicio principal.

\subsection*{Pausa (10 min)}
Recoger una pregunta anónima o un punto de confusión. Elegir una pregunta que permita distinguir síntoma, causa y propiedad.

\subsection*{Bloque 4 - caso trabajado con falla (40 min)}
Aplicar el mecanismo al caso conductor. Introducir una falla a mitad de ejecución y detenerse en cada transición: quién sabe qué, qué se persiste, qué garantía sigue vigente y qué observa el cliente.

\subsection*{Bloque 5 - taller de decisión (25 min)}
Los grupos aplican P-F-G-S-T a su Trabajo Final. Deben producir un ADR breve con alternativa elegida, alternativa descartada y evidencia pendiente. La evidencia de esta clase es presentación de 12-15 minutos, prototipo, preguntas y evaluación con rúbrica.

\subsection*{Bloque 6 - cierre y exit ticket (15 min)}
Cada estudiante responde: propiedad principal, supuesto decisivo, costo aceptado y condición bajo la cual cambiaría de diseño. Usar las respuestas para iniciar la clase siguiente.

\section{Qué explicar, por qué y cómo verificarlo}
\begin{itemize}
\item Explicar el problema antes del producto, porque reconocer la familia de problema es el resultado cognitivo central.
\item Explicar el invariante, porque permite evaluar configuraciones y no solo repetir un flujo nominal.
\item Explicar el límite, porque presentación descriptiva, respuesta evasiva, límite no reconocido y rúbrica inconsistente puede invalidar supuestos o reducir progreso.
\item Verificar comprensión mediante presentación de 12-15 minutos, prototipo, preguntas y evaluación con rúbrica, no mediante una definición aislada.
\end{itemize}

\section{Preguntas socráticas}
\begin{itemize}
\item ¿Qué sabe cada actor en este punto de la ejecución?
\item ¿Qué propiedad sigue siendo cierta si la respuesta nunca llega?
\item ¿Qué parte es safety y cuál es liveness?
\item ¿Qué alternativa más simple resolvería el problema si relajamos una garantía?
\item ¿Qué observación refutaría nuestra hipótesis de diseño?
\end{itemize}

\section{Errores previsibles y estrategia de corrección}
\begin{itemize}
\item Si el alumno responde con un nombre de tecnología, pedir actores, mensajes y estados.
\item Si confunde timeout con falla, construir dos ejecuciones indistinguibles.
\item Si promete todas las garantías, pedir comportamiento bajo partición o saturación.
\item Si mide solo rendimiento, preguntar cómo evidencia la propiedad semántica.
\item Si la solución es excesiva, pedir la alternativa mínima y el costo marginal de la garantía fuerte.
\end{itemize}

\section{Criterios de evaluación formativa}
Una respuesta excelente declara supuestos, recorre una ejecución, nombra con precisión la garantía, reconoce el costo y propone evidencia reproducible. Una respuesta suficiente identifica el mecanismo y el trade-off principal. Una respuesta insuficiente usa slogans, omite fallas o confunde producto con propiedad.

\section{Preparación docente y bibliografía}
Lectura previa recomendada: Rúbrica del curso, ADRs del proyecto y bibliografía técnica utilizada por cada equipo. Preparar un contraejemplo adicional y una variante del caso en la que cambie el costo del error; esto evita convertir la clase en una lectura de diapositivas.

\end{document}
